\begin{textsample}{POS Dim 1 – am – Score 53.00 – am/am\_ef\_41.txt}  \label{ex:f1_pos_224}
Prezado leitor, é com grande satisfação que lhe convidamos a conhecer o Currículo Amazonense para o ensino da Língua \textbf{Espanhola}.
Para \textbf{tanto}, chamamos a atenção para algumas considerações.
Pensar sobre o Amazonas, naturalmente, perpassa por considerar a diversidade e os encontros característicos da região amazônica e, de \textbf{maneira} ampliada, do Brasil.
Além da diversidade biológica, nesse grandioso estado brasileiro convivem indivíduos componentes de distintos povos, carregando em suas identidades seus aspectos culturais, dos quais, neste momento, destacamos as línguas.
Diferentemente do que pensam algumas pessoas, o Brasil não é um país onde se fala exclusivamente um único idioma, o português.
Na realidade, tomando uma perspectiva mais realista, somos um país multilíngue mesmo antes dos primórdios de sua existência enquanto nação.
Desde o reconhecimento da Língua Brasileira de Sinais — \textbf{LIBRAS}, podemos considerar o Brasil um país com dois idiomas oficiais : português e \textbf{LIBRAS}.
A essa realidade \textbf{somam} -se outras, que apresentam dezenas de línguas trazidas por imigrantes que chegaram e aqui se fixaram : alemão, coreano, francês, \textbf{espanhol}, \textbf{inglês}, italiano, japonês, polonês, entre outros idiomas.
Segundo Oliveira ( 2003 ), no Brasil se falam mais de 200 línguas, das quais cerca de 180 são indígenas.
Além dessas, apesar de consideradas extintas formalmente, no país estão também presentes as línguas de \textbf{matriz} africana, principalmente no léxico e em práticas sociais.
Isso tudo caracteriza a realidade multicultural e multilinguística brasileira.
O investigador ( Idem, 2015 ) complementa, informando que existem mais de 120 municípios brasileiros fronteiriços com países que falam outros idiomas oficiais ( principalmente o \textbf{espanhol} ), havendo 10 em regiões de tríplice fronteira, como as realidades nas regiões amazonenses do alto Solimões e do alto Rio Negro.
Tomando o contexto amazonense, percebemos que a situação se intensifica.
Segundo Monteiro ( 2010 ), no estado, há mais de 50 línguas indígenas faladas no dia a dia de comunidades em distintas regiões.
A presença dessas línguas é tão significativa que, no município de São Gabriel da Cachoeira, por exemplo, localizado no noroeste do Amazonas, região de tríplice fronteira — Brasil/Colômbia/Venezuela —, vivem mais de 20 etnias indígenas e, por força de leis municipais, se cooficializaram 4 idiomas nativos : baniwa, nheengatú, tukano e yanomamê.
Sobre a \textbf{pluralidade} linguística na região, Heufemann-Barria e Teixeira ( 2017 ) afirmam, por sua vez, que, além dessas línguas indígenas, no Amazonas encontramos outros idiomas autóctones, trazidos por imigrantes : - O \textbf{inglês}, de importância turística, especialmente por se tratar o Amazonas de uma região de reconhecidas relevância e importância internacionais, estado que abriga grande parte dos rios e da Floresta Amazônica, com imenso \textbf{potencial} turístico, e que tem recebido grandes \textbf{eventos} internacionais como a Copa do Mundo de Futebol, em 2014, e os Jogos Olímpicos, em 2016 ; - O francês, de importância histórica, principalmente durante e \textbf{devido} à época dourada da borracha, quando Manaus passou a ser conhecida internacionalmente como a Paris dos Trópicos ; - O japonês, de importância histórica e \textbf{econômica} para a região, dada a grande comunidade nipônica presente no estado, a qual tem consolidado indústrias, escolas, centros culturais etc. ; e, mais recentemente, - O coreano, igualmente economicamente importante, \textbf{devido} à presença crescente da comunidade coreana na região, especialmente na capital, em indústrias, em instituições etc.
Além desses idiomas, Teixeira ( 2014 ) destaca o \textbf{espanhol}, língua oficial das nações vizinhas, presente de forma viva e significativa em diversas partes do Amazonas.
De acordo com o investigador, especialmente nas regiões turísticas como a capital e seu entorno e nas regiões de fronteira, como o alto rio Solimões — Brasil/Colômbia/Peru — e o alto rio Negro — Brasil/Colômbia/Venezuela —, o \textbf{espanhol} é língua de : - comunicação turística, usada por visitantes estrangeiros ; - negociação comercial, usada por empresários estrangeiros que estabelecem comércio no estado e por habitantes hispânicos que comercializam artigos para uso diário ; - atendimento à saúde, usada por profissionais hispânicos que vivem no estado e que oferecem atendimento médico à população amazonense, especialmente no interior, principalmente a partir do programa Mais Médicos ; - comunicação entre amigos, dada a presença marcante de comunidades hispânicas, com destaque às comunidades colombianas, peruanas e venezuelanas na capital e, também, no interior ; - uso familiar, especialmente nas comunidades de imigrantes hispânicos na região ; e, ainda, - língua de instrução, presente em instituições de ensino — Educação Básica e Ensino Superior, em distintos municípios do Estado — capital e interior, e, de forma especial, nas regiões de fronteira, onde há casos nos quais o \textbf{espanhol} substitui o português como língua formal de instrução.
Dessa forma, percebe -se claramente a presença viva e marcante de diversas línguas e, de \textbf{maneira} especial, a do \textbf{espanhol} no contexto linguístico amazonense.
Aqui, o \textbf{espanhol} é um idioma \textbf{relevante}, que assume distintas funções sociais.
Tomando o contexto internacional, a convergência de interesses de países em que se falam oficialmente as línguas portuguesa e \textbf{espanhola} tem refletido em seu uso nas negociações e nas relações internacionais, envolvendo instituições, países, regiões e blocos onde esses idiomas são falados : Mercado Comum do Sul, Comunidade dos Países de Língua Portuguesa — CPLP, União Latina — UL, Organização dos Estados Ibero-Americanos — OEI, Zona de Paz e Cooperação do Atlântico Sul — ZPCAS, União Africana — UA, Associação Latino-Americana de Integração — ALADI, Banco Interamericano de Desenvolvimento — BID, Comunidade de Estados Latino-Americanos e Caribenhos — CELAC, Grupo do Rio, Organização dos Estados Americanos — OEA, União Europeia — UE, União de Nações Sul-Americanas — Unasul, Organização do Tratado de Cooperação Amazônica — OTCA, entre outros.
De acordo com o sociolinguista \textbf{espanhol} Moreno Fernández ( 2007 ), o \textbf{espanhol} é um idioma que possui reconhecimento internacional e vantagens para ser aprendido, principalmente se considerarmos o fato de que se trata de um idioma neolatino, como o português, apresentando diversas semelhanças com a língua oficial nacional brasileira — o português.
São algumas vantagens apresentadas pelo \textbf{estudioso} \textbf{espanhol} : - língua homogênea, possuindo pequeno índice de fragmentação ; - idioma de cultura reconhecido internacionalmente, presente em obras de \textbf{autores} e artistas consagrados como Miguel de Cervantes, García Lorca, Quevedo, Pablo Neruda, García Márquez, Vargas Llosa, Rubén Darío, Borges, Salvador Dalí, Botero, Frida Kahlo, entre outros ; - língua oficial — em alguns casos, cooficial — em cerca de 30 países, espalhados pelos continentes americano, europeu, africano e asiático, presente ainda em diversos organismos internacionais, como os mencionados anteriormente ; - língua em expansão pelo mundo em número de falantes.
Sobre essa última vantagem, o Instituto Cervantes, em relatório elaborado sobre a evolução do \textbf{espanhol} no mundo ( 2016 ), afirma que, à época da investigação, o \textbf{espanhol} já era falado por mais de 567 milhões de pessoas — quase 8 \% da população mundial, sendo mais de 472 milhões falantes nativos.
Afirma ainda que, das línguas mais faladas no mundo — chinês, \textbf{inglês}, \textbf{espanhol}, hindi e árabe —, por questões demográficas, enquanto o número de falantes de chinês e de \textbf{inglês} está diminuindo, o número de falantes das demais línguas está crescendo.
No Brasil, apesar de não ser idioma oficial, o \textbf{espanhol} é falado por mais de 460 mil pessoas, revelando sua importância para o contexto nacional.
Essa presença é tão significativa no Amazonas que, de acordo com Santos e Teixeira ( 2016 ), desde a década de 80, o ensino do \textbf{Espanhol} está presente no estado, principalmente a partir da fundação da Associação de Professores de \textbf{Espanhol} do Estado do Amazonas — APE-AM, em 1989, \textbf{ator} social que tem envidado esforços desde sua fundação em prol da \textbf{pluralidade} linguística característica do Amazonas, trabalhando pelo hispanismo e pela consolidação do ensino do \textbf{Espanhol} na região.
Além da APE-AM, os investigadores destacam as ações da Universidade Federal do Amazonas — UFAM, \textbf{ator} social que tem envidado esforços para proporcionar formação inicial e continuada de qualidade a profissionais no Amazonas, atendendo a \textbf{demandas} características da comunidade amazonense, considerando sua \textbf{pluralidade}.
Assim, desde 2003, a Ufam tem formado professores de \textbf{Espanhol} para o \textbf{mercado} laboral amazonense, inicialmente na capital e, a partir de 2005, no interior.
Gostaríamos ainda de destacar \textbf{relevantes} ações de outro \textbf{ator} social no Amazonas, a Secretaria de Educação e Qualidade do Ensino do Amazonas — SEDUC/AM, em prol do ensino de \textbf{Espanhol} no estado e da formação continuada de docentes.
Em 2008, foi firmado convênio da SEDUC/AM com a APE-AM, com a UFAM e com a Embaixada da Espanha no Brasil e, mais recentemente, com o Instituto Federal de Educação, Ciência e Tecnologia do Amazonas — IFAM, a fim de realizar formação continuada aos professores de \textbf{Espanhol} no Amazonas, por meio de cursos de atualização de professores.
A \textbf{edição} de 2018, organizada pelo Departamento de Políticas e Programas Educacionais e pelo Centro de Formação Educacional Padre José Anchieta da SEDUC/AM, com o apoio dos parceiros supra, está programada para ocorrer entre os dias 21 e 25/05/2018.
Essa ação tem sido fundamental para o fortalecimento da formação continuada, com ensino de qualidade cada vez melhor aos \textbf{alunos}.
Percebe -se claramente a importância do \textbf{espanhol} para o Amazonas, dado o envolvimento de distintos \textbf{atores} sociais regionais ( SEDUC/AM, APE-AM, UFAM, IFAM ), bem como de entes internacionais.
Outra \textbf{relevante} ação que envolveu distintos \textbf{atores} sociais se deu no âmbito da formação docente, indo ao encontro da realidade plural e das necessidades do estado.
Assim, juntaram -se SEDUC/AM, UFAM, Governo Federal e Secretarias Municipais de Educação ( Barcelos, Santa Izabel do Rio Negro e São Gabriel da Cachoeira ) em prol da formação de professores de \textbf{Espanhol} por meio do Plano Nacional de Formação de Professores da Educação Básica — PARFOR.
Por meio dessa parceria de sucesso, mais de 70 novos professores de \textbf{Espanhol} foram formados no interior do Amazonas, para atuarem nas regiões do \textbf{médio} e do alto rio Negro.
Heufemann-Barría e Teixeira ( 2017, p. 139 ) afirmam ainda que : > No que concerne à UFAM, segundo dados da Pró-Reitoria de Ensino de Graduação — PROEG, no âmbito do PARFOR, foram oferecidos distintos Cursos de Licenciatura, entre os quais : Artes Plásticas, Arte — Música, \textbf{Ciências} Biológicas, \textbf{Ciências} Naturais, Educação Física, Física, Geografia, História, Letras — Língua e Literatura \textbf{Espanhola}, Letras — Língua e Literatura \textbf{Inglesa}, Letras — Língua e Literatura Portuguesa, Licenciaturas Indígenas : Políticas Educacionais e Desenvolvimento Sustentável, Matemática, Pedagogia, Química e Sociologia, atendendo a \textbf{demandas} da capital amazonense e de municípios do interior do Estado.
Realmente foi uma significativa e acertada ação conjunta em prol da formação inicial e continuada de professores no Amazonas, contemplando distintas áreas, entre as quais a de \textbf{Espanhol}.
Para Santos e Teixeira ( op. cit., p. 163 ) : > Essas ações impulsionaram outras.
Projetos de ensino, pesquisa e extensão sobre o hispanismo na capital e no interior do Amazonas têm sido desenvolvidos pela UFAM.
Destacamos uma série de \textbf{eventos} realizados, “ \textbf{Seminários} de Hispanistas do Alto Rio Negro ”, [... ] \textbf{fomentando} o hispanismo em uma região marcada pela \textbf{pluralidade} sociolinguística e cultural, considerada a mais plurilíngue do continente americano, com um sistema ecolinguístico composto por mais de vinte diferentes línguas indígenas, provenientes de cinco troncos distintos : tupi ( nheengatú ), tukano oriental ( tukano, tuyuka, desano, wanano, piratapuya etc. ), aruak ( baniwa, kuripa-ko, tariano, werekena ), makú ( nadeb, daw, yahup, hupda ) e yanomami, além de duas línguas da família românica, o português e o \textbf{espanhol}.
Situação parecida se \textbf{verifica} na região do alto rio Solimões, onde a opção pelos sistemas de ensino como língua estrangeira tem sido o \textbf{Espanhol} ( COELHO e TEIXEIRA, 2014, p. 33 ).
Segundo Coelho ( 2014, p. 31 ), além disso, “ [... ] el Español está presente en las escuelas del Umariaçu, Cordeirinho y Filadélfia, y además podemos afirmar que el Español forma parte del cotidiano de estos indígenas, ubicados en zona de triple frontera. ” Percebemos, mais uma vez, as funções sociais que o \textbf{Espanhol} exerce no Amazonas, também em meio indígena, principalmente nas regiões de fronteira.
Sobre o contexto da região, Guerreiro e Teixeira ( 2017 ) \textbf{reiteram} a escolha do \textbf{Espanhol} como língua de ensino nas escolas dos sistemas estadual e municipais e destacam, ainda, que a decisão se sustenta \textbf{devido} à proximidade com os países hispânicos vizinhos e, consequentemente, ao fluxo intenso de hispanofalantes na região.
No que concerne ao ensino do idioma no Amazonas, segundo Guerreiro ( 2017 ), há registros que comprovam que o \textbf{Espanhol} tem sido ensinado desde o ano de 1997 no alto rio Solimões — Primeira Série do Segundo Grau, na Escola Estadual Imaculada Conceição ( sistema estadual de ensino ), em Benjamin Constant.
O investigador complementa informando que também há registros do ensino de \textbf{Espanhol} na rede municipal, datados do ano de 1998 — Quinta Série do Primeiro Grau, na Escola Municipal Professora Graziela Corrêa de Oliveira ( sistema municipal de ensino ), também em Benjamin Constant.
São dados importantes que revelam e ratificam a relevância do \textbf{espanhol} para o Amazonas e que seu ensino no estado é de fato histórico e deve ser mantido, \textbf{tanto} no sistema estadual, como nos das demais esferas.
Considerando a \textbf{pluralidade} linguística do Amazonas, destacamos ainda outra ação significativamente \textbf{relevante} da SEDUC/AM, a implementação das escolas bilíngues — Português/Japonês, Português/Francês e, mais recentemente, Português/Espanhol ( com o \textbf{relevante} apoio do Consulado Geral da Colômbia, da UFAM e da APE-AM ), e Português/Inglês.
Entendemos, dessa forma, que o poder público estadual amazonense está sensível à realidade do estado e a suas características e, inovando, lança e amplia projeto educacional de grande envergadura.
Destacamos ainda a incorporação de um novo \textbf{ator} social em prol do hispanismo e do ensino de \textbf{Espanhol} no Amazonas, o Consulado Geral da Colômbia.
Entre as ações do Consulado da Colômbia, destacamos o fomento a projetos como o projeto de ensino do \textbf{Espanhol} e da cultura colombiana na Escola Estadual Padre João Badalotti, em Barcelos, no interior do estado, e o projeto Manaus Internacional : integrando culturas por meio da língua \textbf{espanhola}, realizado com o apoio da Secretaria de Educação de Manaus — SEMED/Manaus, na Escola Municipal Raimundo Theodoro, ação de planejamento linguístico para o ensino de \textbf{Espanhol} a crianças.
O envolvimento do governo municipal manauara, reconhecendo a importância do Ensino do \textbf{Espanhol} na capital amazonense, levou a outra ação em prol da difusão do Hispanismo em Manaus, através da Secretaria Municipal de Cultura — ManausCult.
Por cerca de um mês, em distintos pontos culturais da capital amazonense, em 2017, realizou -se o \textbf{evento} Don Quijote — Sonhando um sonho impossível, em comemoração ao Dia Internacional do Livro.
Professores, acadêmicos, \textbf{alunos} de escolas de ensino fundamental e de ensino \textbf{médio} e a sociedade manauara em geral tiveram a oportunidade de conhecer um pouco mais sobre a cultura \textbf{espanhola}, sobre o escritor destacado e sua obra e, principalmente, sobre a língua \textbf{espanhola}.
Mais uma vez, percebemos a relevância do \textbf{espanhol} no Amazonas, dado o envolvimento cada vez maior de \textbf{atores} sociais, entes políticos engajados na difusão do hispanismo e na consolidação do ensino do idioma — SEDUC/AM, APE-AM, UFAM, IFAM, Secretarias Municipais de Ensino ( Manaus, Benjamin Constant, Barcelos, São Gabriel da Cachoeira, Santa Izabel do Rio Negro ), Secretaria Municipal de Cultura de Manaus e, ainda, de \textbf{atores} sociais internacionais como a Embaixada da Espanha e o Consulado Geral da Colômbia.
Nessa mesma esteira, destacamos ainda parceria bem-sucedida entre a UFAM, a APE-AM e a SEDUC/AM, resultando na publicação do livro Ensinando \textbf{Espanhol} no Amazonas : experiências, conquistas e perspectivas ( TEIXEIRA et al., 2017 ), obra que apresenta relatos de experiências desenvolvidas em escolas amazonenses, com destaque para escolas estaduais na capital e no interior, por professores de \textbf{Espanhol} associados à APE-AM.
Essa parceria \textbf{segue} frutífera, de forma que o segundo volume está em fase de organização e, em breve, \textbf{esperamos} que possa ser \textbf{publicado} em \textbf{evento} com a participação dos parceiros mencionados.
Além do \textbf{exposto}, enfatizamos que o ensino do \textbf{Espanhol} no Amazonas tem sido destaque em diversos \textbf{eventos} regionais, nacionais e internacionais, havendo destaque para as parcerias e para as ações mencionadas : - Encontro Internacional da Hispanidade ( Boa Vista/RR, 2011 ) ; - Simpósio Internacional de Letras Neolatinas ( Rio de Janeiro/RJ, 2011 ) ; - Simpósio Linguagens e Identidades da/na Amazônia Sul-Ocidental e Colóquio Internacional “ As Amazônias, as Áfricas e as Áfricas na Pan-Amazônia ” ( Rio Branco/AC, 2012 ) ; - Congresso Internacional da Faculdade de Letras da Universidade Federal do Rio de Janeiro ( Rio de Janeiro/RJ, 2013 ) ; - Encontro Internacional da Hispanidade ( Boa Vista/RR, 2014 ) ; - Fórum Internacional de Pedagogia ( Parintins/AM, 2015 ) ; - Congresso Brasileiro de Professores de \textbf{Espanhol} ( São Carlos/SP, 2015 ) ; - Encuentro de Hispanohablantes en Manacapuru, com a participação dos Consulados Gerais da Colômbia, de Cuba e da Venezuela ( Manacapuru/AM, 2016 ) ; - Congresso Internacional de Professores de Línguas Oficiais do Mercosul e Encontro de Associações de Professores de Línguas Oficiais do Mercosul ( Florianópolis/SC, 2016 ) ; - Encuentro Internacional de Español como Lengua Extranjera : enseñanza, aprendizaje y evaluación ( Bogotá/Colômbia, 2016 ) ; - Mostra de Cinema Latino-Americano de Manacapuru ( Manacapuru/AM, 2017 ) ; - Congresso Brasileiro de Professores de \textbf{Espanhol} ( Belém/PA, 2017 ) ; - Congresso Internacional da Associação de Linguística e Filologia da América Latina ( Bogotá/Colômbia, 2017 ) ; - Congresso Amazônico de Professores de \textbf{Espanhol} ( Macapá/AP, 2018 ) ; - Congresso Nordestino de Professores de \textbf{Espanhol} ( a realizar -se em Natal/RN, 2018 ) ; - Congresso Brasileiro de Hispanistas ( a realizar -se em Aracaju/SE, 2018 ), entre outros.
A relevância e o reconhecimento têm sido tamanhos que o Amazonas foi escolhido, pela primeira vez, para sediar, em 2019, o próximo Congresso Brasileiro de Professores de \textbf{Espanhol}. \textbf{Esperamos} contar com a parceria dos entes sociais envolvidos para o sucesso nessa empreitada!
No que concerne à legislação para o ensino do \textbf{Espanhol}, além de amparo federal, por meio da Lei 13.415/2017 ( BRASIL, 2017 ), que prevê a possibilidade do ensino de um segundo idioma estrangeiro, preferencialmente o \textbf{Espanhol}, no que tange ao âmbito estadual, no Amazonas, o ensino do \textbf{Espanhol} está amparado por meio da Lei 152/13 ( AMAZONAS, 2013 ), que prevê a oferta do idioma neolatino nas escolas amazonenses, conforme o disposto a seguir : > Art. 1º — Fica assegurada a oferta obrigatória da \textbf{disciplina} referente à língua \textbf{espanhola}, nas redes pública e privada do ensino \textbf{médio}, no ato da matrícula dos \textbf{alunos}. > \textbf{Parágrafo} único.
Considera -se oferta obrigatória aquela que se registra mediante manifestação descrita, \textbf{impressa} ou digitada do próprio \textbf{aluno} ou de seu responsável. > Art. 2º — O exercício da atividade de professor de ensino de língua \textbf{espanhola}, no Estado do Amazonas, nas redes pública e privada, é direito exclusivo dos professores formados em curso superior de Letras-Língua \textbf{Espanhola} com licenciatura plena. > Art. 3º — O descumprimento ao disposto na presente lei constitui improbidade administrativa nos termos do Art. 11, I, da Lei Federal n. 8429/1992. > Art. 4º — Esta lei entrará em vigor na data de sua publicação, revogando -se as disposições em contrário.
Para fundamentar o componente de Língua \textbf{Espanhola} no Referencial Curricular Amazonense, foram consideradas as características plurais do Estado e a grande relevância e função social assumidas pelo \textbf{Espanhol} no Amazonas.
Com base em aportes de documentos históricos como os \textbf{Parâmetros} Curriculares Nacionais, as Diretrizes Curriculares Nacionais e, recentemente, a Base Nacional Curricular Comum, aliados ao amparo legal para o ensino do idioma no estado mencionado anteriormente, a construção do documento foi coletiva, havendo a participação de \textbf{colaboradores}, hispanistas de distintos setores da sociedade — Universidade Federal do Amazonas, Seduc, Semed, Associação de Professores de \textbf{Espanhol} do Amazonas, entre outros —, monitorados pela comissão ProBNCC.
A perspectiva de trabalho orienta o componente de \textbf{Espanhol} a partir de cinco eixos.
Oralidade : abrange o uso oral da língua, evidenciando a \textbf{compreensão}, a produção e a interação oral, estimulando o \textbf{aluno} a se envolver em práticas de linguagem oral essencial com contato face a face ( \textbf{debates}, \textbf{entrevistas}, entre outros ), passando a conhecer os “ modos particulares de falar a língua ”.
Alguns aspectos \textbf{relevantes}, como entonação e ritmo, ao serem articulados em consonância com as estratégias de \textbf{compreensão}, são de \textbf{suma} importância para a exploração das práticas de linguagem em situações de uso oral da língua.
Nessas práticas, que articulam não só o aspecto verbal, mas também o visual, o sonoro, o gestual e o tátil, os estudantes terão oportunidades de vivência e de reflexão sobre os usos orais/oralizados do \textbf{Espanhol}.
Leitura : compreende as práticas de linguagem oriundas do entendimento do processo de leitura interativo, quer seja entre leitor-texto-autor, quer seja entre leituras de mundo ; o leitor deve se apoiar na \textbf{compreensão} e na interpretação dos gêneros escritos em língua \textbf{espanhola}, que circulam nos diversos campos e esferas da sociedade.
Nesse eixo, o \textbf{aluno} desenvolve as práticas de leitura necessárias para reconhecer tipologias e gêneros textuais/discursivos, além de \textbf{aprimorar} o senso \textbf{crítico} em percurso criativo e autônomo de aprendizagem da língua.
Além disso, as práticas leitoras em língua \textbf{espanhola} compreendem possibilidades variadas de contextos de uso das linguagens para pesquisa e ampliação de conhecimentos de temáticas significativas para os estudantes, com trabalhos de natureza interdisciplinar e/ou fruição estética de gêneros como poemas, peças de teatro etc.
Escrita : o eixo da escrita considera dois aspectos do ato de escrever — natureza processual e colaborativa e o escrever como prática social —, propondo aos \textbf{alunos} a oportunidade de agir com protagonismo em uma escrita autoral, que se inicia com textos que utilizam poucos recursos verbais ( mensagens, tirinhas, fotolegendas, adivinhas, entre outros ) e se desenvolve para textos mais elaborados ( autobiografias, esquetes, notícias, relatos de opinião, chat, fôlderes, entre outros ), nos quais são utilizados recursos linguístico-discursivos variados, em movimentos coletivos e individuais de planejamento, de produção, de revisão e de reescritura.
Conhecimentos linguísticos : fundamenta -se pelas práticas de uso a serviço das habilidades de leitura, de escrita e de oralidade.
O estudo da gramática tem como objetivo fazer com que o \textbf{aluno} compreenda o funcionamento sistêmico e social da língua \textbf{espanhola}.
Desenvolvem -se noções de variedade padrão e o respeito às demais variedades linguísticas do idioma, explorando suas semelhanças e diferenças e, de modo contrastivo, suas semelhanças e diferenças com a língua portuguesa e com outros idiomas.
Essas noções constituem um exercício metalinguístico efetivo e o reconhecimento de características do próprio idioma do \textbf{aluno}.
Dimensões \textbf{interculturais} : ressalta o conceito do \textbf{espanhol} como língua internacional, analisando os diferentes papéis do idioma no mundo, seus valores, seu alcance e seus efeitos nas relações entre diferentes pessoas e povos, \textbf{tanto} na sociedade \textbf{contemporânea} quanto em uma perspectiva histórica.
Surge da \textbf{compreensão} de que as culturas passam continuamente por um processo de interação e de (re)construção.
Os eixos devem ser trabalhados de forma simultânea, fazendo com que o processo de ensino e de aprendizagem da língua \textbf{espanhola} seja híbrido, polifônico e multimodal, ressaltando que os mesmos se compõem de unidades temáticas, as quais, em sua grande maioria, repetem -se, ampliando as habilidades a elas correspondentes.
Cada unidade temática possui objetos de conhecimentos e habilidades a serem trabalhados no \textbf{decorrer} do ano.
Sendo assim, o componente de língua \textbf{espanhola} está organizado da seguinte forma : eixo e unidade temática, \textbf{competências}, habilidades, objeto de conhecimento e detalhamento do objeto de conhecimento.
Na \textbf{seção} de detalhamento do objeto de conhecimento, encontra -se um espaço com sugestão de vocabulário e outros aspectos a serem trabalhados de acordo com a habilidade.
As \textbf{competências} específicas do componente abaixo relacionadas articulam -se às \textbf{competências} gerais da Base Nacional Curricular Comum ( BNCC ).
Dessa forma, por todo o \textbf{exposto}, encorajamos que se aproprie deste documento, observando a possibilidade de realizar um trabalho diferenciado, atentando especialmente às \textbf{indicações} no campo “ detalhamento do objeto de conhecimento ”, que \textbf{sinaliza} conceitos que perpassam por mais de um componente, marca de \textbf{interdisciplinaridade}, o que facilita o planejamento de suas \textbf{aulas}, oportunizando uma aprendizagem significativa dos \textbf{alunos}.
Texto de \textbf{autoria} dos Profes.
Dres.
Elsa Otilia Heufemann Barria e Wagner Barros Teixeira — Ufam, \textbf{baseado} em capítulo de livro \textbf{publicado} pela Associação de Professores de \textbf{Espanhol} do Amazonas.
Lei 145/02 ( SÃO GABRIEL DA CACHOEIRA, 2002 ).
Lei 0084/17 ( SÃO GABRIEL DA CACHOEIRA, 2017 ).
Programa do Governo Federal que objetiva melhorar a atenção à saúde no país com ações que levem profissionais a regiões onde existe escassez e, ainda, ampliar e construir unidades básicas de saúde, além de aumentar o número de vagas e de melhorar a formação superior na área da saúde ( cf. maismedicos.gov.br ).
Atualmente corresponde ao primeiro ano do ensino \textbf{médio}.
Atualmente corresponde ao sexto ano do ensino fundamental.

% matched lemmas: aluno, aprimorar, ator, aula, autor, autoria, basear, ciência, colaborador, competência, compreensão, contemporâneo, crítico, debate, decorrer, demanda, devido, disciplina, econômico, edição, entrevista, espanhol, esperar, estudioso, evento, exposto, fomentar, imprimir, indicação, inglês, intercultural, interdisciplinaridade, libra, maneira, matriz, mercado, médio, parágrafo, parâmetro, pluralidade, potencial, publicar, reiterar, relevante, segar, seminário, seção, sinalizar, somar, sumo, tanto, verificar
\end{textsample}
